\documentclass[conference]{IEEEtran}

\usepackage[colorlinks,urlcolor=blue,linkcolor=blue,citecolor=blue]{hyperref}
\usepackage{amsmath}
\usepackage{amssymb}
\usepackage{graphicx}

% Reduce title font size (IEEEtran default is \LARGE ~17pt; drop to 13pt)
% and tighten the gap between title block and abstract
\makeatletter
\renewcommand{\@maketitle}{%
  \newpage\null\vskip 0.2em%
  \begin{center}%
    {\fontsize{13}{15}\bfseries\selectfont \@title \par}%
    \vskip 0.75em%
    {\lineskip .5em\begin{tabular}[t]{c}\@author\end{tabular}\par}%
    \vskip 0.4em%
  \end{center}\par\vskip 0.4em}
\makeatother

\begin{document}

\title{LoRA as Declarative Adaptation for Multimodal Clinical Question Answering in a Virtual Hospital Environment}

\author{
  \IEEEauthorblockN{Sandrine de Ribaupierre\IEEEauthorrefmark{1}\IEEEauthorrefmark{2},\;
                     Roy Eagleson\IEEEauthorrefmark{2}\IEEEauthorrefmark{3},\;
                     Anthony Barros\IEEEauthorrefmark{3},\;
                     Kelsey Kloosterman\IEEEauthorrefmark{3}}
  \IEEEauthorblockA{\IEEEauthorrefmark{1}Western University CNS\quad
                     \IEEEauthorrefmark{2}Centre for Brain and Mind, Western University\\
                     \IEEEauthorrefmark{3}Faculty of Engineering, Dept.\ of AI and Software Engineering, Western University}
}

\maketitle

\begin{abstract}
Pediatric patients in hospital environments require age-appropriate, institutionally grounded information that general-purpose AI systems are not equipped to provide. We present a LoRA fine-tuning framework for clinical question-answering within a virtual hospital environment, targeting two pediatric age groups (5--11 and 12--18). Using Low-Rank Adaptation (LoRA)~\cite{hu2022lora} on 4-bit quantized Llama~3.2 Instruct models (1B and 3B), we train Standard ($r{=}8$, 16 layers) and Fast ($r{=}4$, 8 layers) adapter configurations, producing eight training runs across model sizes and age groups. Fine-tuned adapters substantially outperform zero-shot base models on age-appropriate readability~\cite{flesch1948,gunning1952,senter1967,mclaughlin1969,kincaid1974,kincaid1975,coleman1975} (Flesch-Kincaid Grade~$\leq$7.0 pass rates: 72.84\% vs.\ 12.14\%) and inter-role style separation (0.89--0.94 vs.\ 0.66--0.70). We observe a configuration-ordering crossover between 1B and 3B models: the Standard adapter outperforms on 3B (84\% vs.\ 72\%), while Fast outperforms on 1B (82\% vs.\ 76\%), consistent across all five readability metrics. Only Fast-1B meets the 1.0\,s VR latency target. These results demonstrate that compact LoRA adaptation on consumer hardware is feasible for pediatric clinical NLP, but adapter configurations do not transfer reliably across model scales~\cite{biderman2024}.
\end{abstract}

\begin{IEEEkeywords}
Clinical Q\&A systems, Pediatric NLP, LoRA, Virtual Reality
\end{IEEEkeywords}

\section{Introduction}

We present results situated within a virtual hospital framework previously developed by our group to support interactive, clinically grounded communication among pediatric patients, caregivers, and clinical teams. In this work, we extend the platform with a constrained multimodal question-answering capability integrating large language models (LLMs) and vision-language models (VLMs) to enable safe, naturalistic dialogue aligned with institutional guidance and clinical workflows. Our primary objective is to determine whether Low-Rank Adaptation (LoRA)~\cite{hu2022lora} can produce quantifiably age-appropriate pediatric question answering~\cite{flesch1948,gunning1952,senter1967,mclaughlin1969,kincaid1974,kincaid1975,coleman1975} with real-time local deployment on commodity graphical workstations, and how optimal adapter configurations depend on model scale.

\section{Materials and Methods}

Using Low-Rank Adaptation (LoRA)~\cite{hu2022lora} on 4-bit quantized Llama~3.2 Instruct models (1B and 3B)~\cite{dettmers2023}, we locally train two adapter configurations: Standard ($r{=}8$, 16 layers) and Fast ($r{=}4$, 8 layers), across both model sizes and age groups for 8 training runs total. Training data consists of 500 synthetic examples per adapter, generated by Claude Opus~4.6 using Victoria Hospital's public patient materials as grounding artifacts, following a LIMA-inspired curation approach~\cite{zhou2023lima}.

We conceptualize LoRA as a form of \textit{declarative adaptation}: structured institutional knowledge is encoded as modular behavioral constraints rather than procedural rules. Hospital guidelines, care pathways, and patient information resources define a specification space of acceptable responses, while the adapted model realizes that specification during interaction. We further introduce a five-factor corpus quality framework---correctness, consistency, completeness, conclusiveness, and relevance---used both to construct controlled synthetic training corpora and to evaluate downstream responses.

\section{Results}

Fine-tuned adapters substantially outperform zero-shot base models on age-appropriate readability~\cite{flesch1948,gunning1952,senter1967,mclaughlin1969,kincaid1974,kincaid1975,coleman1975} (Flesch-Kincaid Grade~7.0 pass rates: 72.84\% vs.\ 12.14\%) and inter-role style separation (0.89--0.94 vs.\ 0.66--0.70). A notable empirical finding is a \textbf{configuration-ordering crossover} across model scales~\cite{biderman2024}. Under a Standard configuration, 3B models outperform 1B on readability (84\% vs.\ 72\% pass rate); under a Fast configuration this ordering reverses, with Fast-1B outperforming Fast-3B (82\% vs.\ 76\%). This crossover is directionally consistent across all five readability measures examined. Latency results indicate that only Fast-1B satisfies the sub-second response target required for immersive VR interaction.

\section{Discussion and Conclusion}

We evaluated the resulting system in a medical visual question-answering setting, where users pose natural language queries regarding procedures, imaging workflows, expected duration, comfort measures, and care safeguards. Performance was assessed through automated readability metrics~\cite{flesch1948,gunning1952,senter1967,mclaughlin1969,kincaid1974,kincaid1975,coleman1975}, inter-role style separation, latency measurements, and expert inspection. Our findings suggest that compact LoRA adaptation~\cite{hu2022lora} provides a practical pathway for privacy-preserving pediatric clinical NLP on consumer hardware, but adapter configurations do not transfer reliably across model scales~\cite{biderman2024}. Combining synthetic corpus design~\cite{zhou2023lima}, multimodal adaptation, and local deployment offers a scalable and governable route toward clinical communication systems in virtual healthcare environments.

\vspace{4pt}
\noindent\textit{Disclosures:} The authors declare no conflict of interest.

\begin{thebibliography}{11}

\bibitem{flesch1948}
R.~Flesch, ``A new readability yardstick,'' \textit{J. Applied Psychology}, vol.~32, no.~3, pp.~221--233, 1948.

\bibitem{gunning1952}
R.~Gunning, \textit{The Technique of Clear Writing}. New York: McGraw-Hill, 1952, pp.~36--37.

\bibitem{senter1967}
R.~J.~Senter and E.~A.~Smith, ``Automated readability index,'' AMRL-TR-66-20, Wright-Patterson AFB, 1967.

\bibitem{mclaughlin1969}
G.~H.~McLaughlin, ``SMOG grading---a new readability formula,'' \textit{J. Reading}, vol.~22, pp.~639--646, 1969.

\bibitem{kincaid1974}
J.~P.~Kincaid and W.~C.~McDaniel, ``An inexpensive automated way of calculating Flesch Reading Ease scores,'' Patent Disclosure 031~350, U.S. Patent Office, 1974.

\bibitem{kincaid1975}
J.~P.~Kincaid, L.~Fishburne, R.~Rogers, and B.~Chissom, ``Derivation of new readability formulas for Navy enlisted personnel,'' Research Branch Report 8-75, Naval Technical Training Command, 1975.

\bibitem{coleman1975}
M.~Coleman and T.~L.~Liau, ``A computer readability formula designed for machine scoring,'' \textit{J. Applied Psychology}, vol.~60, pp.~283--284, 1975.

\bibitem{hu2022lora}
E.~J.~Hu, Y.~Shen, P.~Wallis, et al., ``LoRA: Low-Rank Adaptation of large language models,'' in \textit{Proc. ICLR}, 2022.

\bibitem{dettmers2023}
T.~Dettmers, A.~Pagnoni, A.~Holtzman, and L.~Zettlemoyer, ``QLoRA: Efficient finetuning of quantized LLMs,'' in \textit{Adv. Neural Inf. Process. Syst.}, 2023.

\bibitem{biderman2024}
D.~Biderman, J.~G.~Ortiz, J.~Portes, et al., ``LoRA learns less and forgets less,'' \textit{Trans. Mach. Learn. Res.}, 2024.

\bibitem{zhou2023lima}
C.~Zhou, P.~Liu, P.~Xu, et al., ``LIMA: Less is more for alignment,'' in \textit{Adv. Neural Inf. Process. Syst.}, 2023.

\end{thebibliography}

\end{document}
